\documentclass[11pt,a4paper]{article}

% ===================== PACKAGES =====================
\usepackage[T1]{fontenc}
\usepackage[utf8]{inputenc}
\usepackage[a4paper,margin=2.5cm]{geometry}
\usepackage{amssymb}
\usepackage[table]{xcolor}
\usepackage{booktabs}
\usepackage{longtable}
\usepackage{array}
\usepackage{tabularx}
\usepackage{enumitem}
\usepackage{titlesec}
\usepackage[hidelinks]{hyperref}

% ===================== STYLING =====================
\definecolor{antamblue}{RGB}{0,51,102}
\definecolor{antamgold}{RGB}{197,148,40}
\definecolor{softgray}{RGB}{240,240,240}

\titleformat{\section}{\color{antamblue}\Large\bfseries}{\thesection.}{0.6em}{}
\titleformat{\subsection}{\color{antamblue}\large\bfseries}{\thesubsection}{0.6em}{}

\renewcommand{\contentsname}{Daftar Isi}
\newcommand{\todo}{\item[$\square$]}
\newcolumntype{Y}{>{\raggedright\arraybackslash}X}

\hypersetup{
  pdftitle={NickelScope - Rencana Kerja},
  pdfauthor={Tim NickelScope}
}

% ===================== TITLE =====================
\title{\color{antamblue}\textbf{NickelScope}\\[0.3em]
\large Rencana Kerja \& Alur Pengerjaan Lengkap\\[0.2em]
\normalsize ANTAM Hackathon 2026 --- Young Mining Innovators\\
\normalsize Tema: Eksplorasi \textbar{} AOI: Pomalaa--Kolaka, Sulawesi Tenggara}
\author{Tim NickelScope}
\date{\today}

\begin{document}
\maketitle
\thispagestyle{empty}

\begin{center}
\fcolorbox{antamgold}{softgray}{\parbox{0.92\textwidth}{\centering
\textbf{Ringkasan:} Sistem AI--GIS yang memetakan probabilitas keterdapatan
nikel laterit dari fusi citra satelit, DEM, dan peta geologi --- menghasilkan
peta target bor prioritas sekaligus layer risiko ESG untuk menekan biaya,
waktu, dan dampak lingkungan eksplorasi.}}
\end{center}

\vspace{0.5em}
\tableofcontents
\newpage

% ============================================================
\section{Gambaran Umum Proyek}
\label{sec:overview}

\begin{description}[leftmargin=3.2cm, style=nextline]
  \item[\textbf{Judul}] NickelScope --- AI-GIS Prospectivity Mapping untuk
        Eksplorasi Nikel Laterit Berbasis Penginderaan Jauh.
  \item[\textbf{Tema}] Eksplorasi (komoditas nikel --- prioritas utama ANTAM).
  \item[\textbf{AOI}] Pomalaa--Kolaka, Sulawesi Tenggara.
        Bounding box: Lat $-4.35^\circ$ s/d $-3.95^\circ$~S;
        Lon $121.50^\circ$ s/d $121.75^\circ$~E ($\sim$28\,km $\times$ 44\,km).
  \item[\textbf{Data}] 100\% publik \& legal (Sentinel-2, DEMNAS/SRTM,
        peta geologi Badan Geologi, konsesi MOMI/MODI).
  \item[\textbf{Output}] WebGIS dashboard interaktif: peta prospektivitas,
        ranking target bor, layer no-go ESG, peta ketidakpastian.
\end{description}

\subsection*{Alur Pipeline Solusi}
\begin{center}\footnotesize
\fbox{Akuisisi Data} $\rightarrow$ \fbox{Pelabelan} $\rightarrow$
\fbox{Feature Engineering} $\rightarrow$ \fbox{Pemodelan ML}\\[0.4em]
$\rightarrow$ \fbox{Validasi \& Uncertainty} $\rightarrow$
\fbox{Peta Prospektivitas} $\rightarrow$ \fbox{Dashboard WebGIS}
\end{center}

% ============================================================
\section{Timeline \& Milestone Lomba}
\label{sec:timeline}

Strategi: \textbf{target utama = proposal PDF tuntas sebelum 10 Juli 2026}.
Bila lolos finalis, prototype diperdalam untuk presentasi (22 Juli) dan
Demo Day (11 Agustus).

\renewcommand{\arraystretch}{1.3}
\begin{longtable}{p{2.7cm} p{3.0cm} p{6.2cm} p{2.9cm}}
\toprule
\rowcolor{softgray}
\textbf{Periode} & \textbf{Milestone Lomba} & \textbf{Aktivitas Internal Tim} & \textbf{Output} \\
\midrule
\endhead
1--15 Jun & Pendaftaran & Daftar tim, finalisasi peran, setup tools (akun GEE, repo) & Tim terdaftar \\
18 Jun & Workshop Teknis & Ikuti workshop; serap arahan panitia tentang data/tools & Catatan teknis \\
18 Jun -- 1 Jul & (Pengerjaan) & Fase 1--3: akuisisi data, pelabelan, feature engineering, model awal & Dataset + model v1 \\
1--9 Jul & (Pengerjaan) & Fase 4--6: validasi, dashboard MVP, tulis proposal PDF & Draf proposal \\
\textbf{10 Jul} & \textbf{Batas Proposal} & Submit PDF proposal final & \textbf{Proposal terkirim} \\
15 Jul & Pengumuman Finalis & (Menunggu hasil) & --- \\
17 Jul & Sprint \& Mentoring & Perdalam prototype, siapkan slide \& video pitch, sesi mentor & Slide + video \\
22 Jul & Presentasi Finalis & Presentasi + tanya jawab di hadapan juri & Pitch deck final \\
\textbf{11 Agu} & \textbf{Demo Day} & Live demo prototype di Kantor ANTAM & \textbf{Demo + awarding} \\
\bottomrule
\caption{Pemetaan jadwal lomba ke rencana kerja internal.}
\end{longtable}

% ============================================================
\section{Pembagian Peran Tim}
\label{sec:roles}

Tim 3--4 orang (kekuatan: Geologi/Tambang + Data/AI). Saran peran
(boleh dirangkap bila tim 3 orang):

\begin{longtable}{p{3.2cm} p{11.5cm}}
\toprule
\rowcolor{softgray}
\textbf{Peran} & \textbf{Tanggung Jawab Utama} \\
\midrule
\endhead
\textbf{Domain (Geologi)} & Interpretasi geologi, host-rock constraint, QA label,
validasi geologis hasil model, framing ESG/Safety. \\
\textbf{ML / Data} & Pipeline data, feature engineering, training model
(RF/XGBoost), spatial cross-validation, peta ketidakpastian. \\
\textbf{GIS / RS} & Google Earth Engine, ekstraksi fitur spasial, pembuatan
peta \& dashboard WebGIS (dapat dirangkap ML/Data). \\
\textbf{Lead / Presenter} & Koordinasi, penulisan proposal, slide deck,
storyline pitch, manajemen deadline. \\
\bottomrule
\end{longtable}

% ============================================================
\section{Alur Pengerjaan per Fase}
\label{sec:phases}

Setiap fase memuat \textbf{Tujuan}, \textbf{Tugas} (checklist),
dan \textbf{Definition of Done (DoD)}.

% ---------- FASE 0 ----------
\subsection{Fase 0 --- Persiapan \& Setup}
\textbf{Tujuan:} Menyiapkan fondasi kerja tim.
\begin{itemize}[leftmargin=1.5em]
  \todo Daftar tim ke panitia (CP: Satriya 0812-1829-1119 / Bella 0813-4061-1919).
  \todo Buat akun Google Earth Engine (gratis) untuk semua anggota.
  \todo Siapkan repo bersama (GitHub) + struktur folder (\texttt{data/}, \texttt{gee/}, \texttt{notebooks/}, \texttt{dashboard/}, \texttt{docs/}).
  \todo Install tools: Python (pandas, scikit-learn, xgboost, geopandas, rasterio), QGIS.
  \todo Sepakati peran \& jadwal mingguan.
\end{itemize}
\textbf{DoD:} Semua anggota punya akses GEE + repo; lingkungan Python siap.

% ---------- FASE 1 ----------
\subsection{Fase 1 --- Akuisisi Data \& Pelabelan}
\textbf{Tujuan:} Menyusun dataset titik berlabel (positif/negatif) yang valid.
\begin{itemize}[leftmargin=1.5em]
  \todo Tetapkan AOI di GEE (\texttt{gee/01\_aoi\_explore.js}) --- \emph{sudah dibuat}.
  \todo Unduh peta geologi 1:250.000 (Badan Geologi) $\rightarrow$ ekstrak unit \textbf{ultramafik} (host-rock mask).
  \todo Ambil polygon konsesi nikel (WIUP/IUP) dari MOMI/MODI ESDM.
  \todo Digitasi footprint tambang dari citra (iron-oxide tinggi + NDVI rendah) $\rightarrow$ \textbf{titik positif}.
  \todo Kumpulkan koordinat okurensi dari publikasi (profil Pomalaa, Pondidaha, dll).
  \todo Generate \textbf{titik negatif}: di litologi non-ultramafik \& area jauh dari okurensi.
  \todo Ekspor menjadi \texttt{labels.csv} (kolom: \texttt{lat, lon, label}).
\end{itemize}
\textbf{DoD:} \texttt{labels.csv} berisi $\geq$ beberapa ratus titik seimbang,
sudah di-QA oleh anggota geologi.

% ---------- FASE 2 ----------
\subsection{Fase 2 --- Feature Engineering}
\textbf{Tujuan:} Menghasilkan stack fitur prediktor berbasis pengetahuan geologi.
\begin{itemize}[leftmargin=1.5em]
  \todo \textbf{Spektral (Sentinel-2):} Iron Oxide Ratio (B4/B2), NDVI, Clay/Ferrous index.
  \todo \textbf{Topografi (DEM):} elevasi, slope, curvature, TWI (Topographic Wetness Index).
  \todo \textbf{Geologi:} jarak ke litologi ultramafik, densitas lineament/struktur.
  \todo Susun semua fitur dalam satu image stack; sampling nilai fitur di tiap titik label.
  \todo Ekspor tabel fitur+label untuk training (\texttt{features.csv}).
\end{itemize}
\textbf{DoD:} \texttt{features.csv} siap; tidak ada nilai kosong/anomali kritis.

% ---------- FASE 3 ----------
\subsection{Fase 3 --- Pemodelan Machine Learning}
\textbf{Tujuan:} Melatih model prospektivitas yang akurat \& dapat dijelaskan.
\begin{itemize}[leftmargin=1.5em]
  \todo Baseline: Weights of Evidence (WoE) --- metode klasik, mudah dijelaskan ke juri.
  \todo Model utama: Random Forest \& XGBoost (bandingkan performa).
  \todo Terapkan host-rock mask sebelum prediksi (buang area non-ultramafik).
  \todo Prediksi probabilitas ke seluruh AOI $\rightarrow$ raster prospektivitas 0--1.
  \todo Feature importance $\rightarrow$ verifikasi konsistensi geologis.
\end{itemize}
\textbf{DoD:} Raster peta prospektivitas + ranking target prioritas dihasilkan.

% ---------- FASE 4 ----------
\subsection{Fase 4 --- Validasi \& Peta Ketidakpastian}
\textbf{Tujuan:} Membuktikan model kredibel \& jujur soal ketidakpastian.
\begin{itemize}[leftmargin=1.5em]
  \todo \textbf{Spatial / block cross-validation} (BUKAN random split) --- hindari kebocoran data spasial.
  \todo Metrik: ROC--AUC, success-rate curve, confusion matrix.
  \todo Peta ketidakpastian (variance/entropy prediksi).
  \todo Validasi geologis: cek apakah zona high-probability masuk akal secara geologi.
\end{itemize}
\textbf{DoD:} Laporan validasi + peta ketidakpastian; metrik terdokumentasi.

% ---------- FASE 5 ----------
\subsection{Fase 5 --- Prototype Dashboard (WebGIS)}
\textbf{Tujuan:} Membungkus hasil jadi prototype interaktif yang ``wow''.
\begin{itemize}[leftmargin=1.5em]
  \todo Bangun dashboard (Streamlit + Folium/Leaflet atau GEE App).
  \todo Layer: peta prospektivitas, ranking target bor, \textbf{no-go ESG}, ketidakpastian.
  \todo Fitur interaktif: klik area $\rightarrow$ tampil skor, koordinat, tingkat keyakinan.
  \todo Data quality report otomatis.
\end{itemize}
\textbf{DoD:} Dashboard dapat di-demo end-to-end secara langsung.

% ---------- FASE 6 ----------
\subsection{Fase 6 --- Penulisan Proposal (Submission Seleksi)}
\textbf{Tujuan:} Proposal PDF sesuai TOR bagian F (deadline 10 Juli).
\begin{itemize}[leftmargin=1.5em]
  \todo Judul \& tema; Problem statement; Cakupan studi kasus.
  \todo Data \& metode; Target prototype; Dampak \& manfaat; Feasibility \& roadmap.
  \todo Sisipkan figur: peta prospektivitas, diagram pipeline, layer ESG.
  \todo Proofread \& cek kesesuaian dengan kriteria penjurian.
\end{itemize}
\textbf{DoD:} PDF proposal final terkirim sebelum 10 Juli 2026.

% ---------- FASE 7 ----------
\subsection{Fase 7 --- Sprint Finalis}
\textbf{Tujuan:} Siapkan materi finalis (bila lolos).
\begin{itemize}[leftmargin=1.5em]
  \todo Slide deck (maks 15 slide).
  \todo Dokumen pendukung: kode/script/desain dashboard.
  \todo Video pitch (maks 2 menit) sesuai storyline.
  \todo Latihan presentasi \& antisipasi pertanyaan juri (lintas disiplin).
\end{itemize}
\textbf{DoD:} Slide + video + kode siap; tim sudah berlatih pitch.

% ---------- FASE 8 ----------
\subsection{Fase 8 --- Demo Day}
\textbf{Tujuan:} Presentasi \& live demo final di Kantor ANTAM (11 Agustus).
\begin{itemize}[leftmargin=1.5em]
  \todo Siapkan demo offline (antisipasi koneksi internet terbatas).
  \todo Bawa analisis manfaat \& rencana implementasi.
  \todo Pegang storyline: hook $\rightarrow$ problem $\rightarrow$ solusi $\rightarrow$ demo $\rightarrow$ impact.
\end{itemize}
\textbf{DoD:} Live demo berjalan mulus; pertanyaan juri terjawab.

% ============================================================
\section{Manajemen Risiko}
\label{sec:risk}

\renewcommand{\arraystretch}{1.3}
\begin{longtable}{p{4.3cm} p{4.0cm} p{6.2cm}}
\toprule
\rowcolor{softgray}
\textbf{Risiko} & \textbf{Dampak} & \textbf{Mitigasi} \\
\midrule
\endhead
Citra Sentinel-2 berawan & Fitur spektral buruk & Gunakan musim kering (Jul--Okt), median composite multi-tanggal. \\
Kurva belajar GEE & Pengerjaan lambat & Mulai dari script starter; bagi tugas; banyak tutorial publik. \\
Label tidak akurat & Model menyesatkan & Host-rock constraint + QA geologi + multi-sumber label. \\
Negatif = belum tereksplorasi & Bias model & Dokumentasikan di peta ketidakpastian; jujur di proposal. \\
Overfitting / data leakage & AUC palsu tinggi & Spatial/block CV, bukan random split. \\
Waktu mepet ke deadline & Proposal tidak tuntas & Prioritaskan MVP; kerja paralel (data \& proposal). \\
\bottomrule
\end{longtable}

% ============================================================
\section{Pemetaan ke Kriteria Penjurian}
\label{sec:scoring}

\renewcommand{\arraystretch}{1.3}
\begin{longtable}{p{5.2cm} p{1.5cm} p{7.8cm}}
\toprule
\rowcolor{softgray}
\textbf{Kriteria} & \textbf{Bobot} & \textbf{Bagaimana NickelScope Menjawab} \\
\midrule
\endhead
Problem statement \& metodologi & 25\% & Masalah nyata komoditas prioritas + metode WoE+ML teruji. \\
Kreativitas \& teknologi & 25\% & Fusi geospasial + ML + feature engineering geologi. \\
Kualitas prototype \& visualisasi & 25\% & WebGIS interaktif, peta visual yang demoable. \\
Dampak \& potensi implementasi & 15\% & Cost saving terukur + roadmap integrasi ANTAM. \\
Presentasi \& tanya jawab & 10\% & Storyline kuat + demo peta yang dramatis. \\
\bottomrule
\end{longtable}

% ============================================================
\section{Checklist Deliverable per Tahap Lomba}
\label{sec:deliverables}

\begin{itemize}[leftmargin=1.5em]
  \todo \textbf{Seleksi (10 Jul):} PDF proposal (masalah, metodologi, studi kasus, target prototype, dampak, feasibility).
  \todo \textbf{Sprint Finalis (17 Jul):} Slide deck ($\leq$15 slide) + dokumen kode/dashboard + video pitch ($\leq$2 menit).
  \todo \textbf{Demo Day (11 Agu):} Live demo prototype + analisis manfaat + rencana implementasi.
\end{itemize}

\vspace{1em}
\begin{center}\itshape
``INNOVATE TODAY, IMPACT TOMORROW.''
\end{center}

\end{document}
